\newpage
\phantomsection
\label{prf:every-element-is-one-or-a-successor}
\LRAProofFor{thm:every-element-is-one-or-a-successor}

\begin{remark*}[Return]
\hyperref[thm:every-element-is-one-or-a-successor]{Return to Theorem}
\end{remark*}

\begin{theorem*}[Every Element Is Either $1$ or a Successor]
Let $(P,S,1)$ be a Peano system. For every $x \in P$, either $x=1$ or there
exists $u \in P$ such that $S(u)=x$.
\end{theorem*}

\begin{proof}
\textbf{Professional Standard Proof.}~\\
Let
\[
D:=\{x\in P:x=1 \text{ or } \exists u\in P\;(S(u)=x)\}.
\]
We show that $D$ is inductive. Since $1=1$, we have $1\in D$.

Now let $x\in D$. By closure of successor, $S(x)\in P$. Moreover $x\in P$,
so $S(x)$ is the successor of an element of $P$. Hence $S(x)\in D$.
Therefore $D$ contains $1$ and is closed under successor. By the induction
principle, $D=P$.

Thus every $x\in P$ lies in $D$, so every $x\in P$ satisfies either $x=1$ or
there exists $u\in P$ such that $S(u)=x$.
\end{proof}

\begin{proof}
\textbf{Detailed Learning Proof.}~\\
Define a subset $D$ of $P$ by putting into $D$ exactly those elements of $P$
that already have the desired form:
\[
D=\{x\in P:x=1 \text{ or } \exists u\in P\;(S(u)=x)\}.
\]
We prove that $D=P$ by induction.

For the base case, $1\in D$ because $1=1$. Thus the distinguished first
element has the required form.

For the successor step, assume $x\in D$. Since $D\subseteq P$, we have
$x\in P$. The Peano closure axiom gives $S(x)\in P$. But then $S(x)$ is
literally a successor, namely the successor of $x\in P$. Therefore
\[
\exists u\in P\;(S(u)=S(x)),
\]
for example by taking $u=x$. Hence $S(x)\in D$.

So $D$ contains $1$ and is closed under successor. By the induction principle
for Peano systems, $D=P$. Since membership in $D$ is exactly the assertion
that an element is either $1$ or a successor, the theorem follows.
\end{proof}

\begin{remark*}[Proof structure]
The proof forms the subset of all elements satisfying the desired dichotomy
and then proves that this subset is inductive. The base case is immediate
from $1=1$. The successor step is immediate because every successor is, by
definition, a successor of an element of $P$.
\end{remark*}

\begin{dependencies}
\begin{itemize}
  \item \hyperref[def:peano-system]{Peano System}
  \item \hyperref[ax:peano-one-in-n]{Distinguished Element Lies in the Underlying Set}
  \item \hyperref[ax:peano-successor-closure]{Closure Under Successor}
  \item \hyperref[thm:induction-principle-for-peano-system]{Induction Principle for a Peano System}
\end{itemize}
\end{dependencies}

